\documentclass[12pt]{ctexart}
\usepackage{amsmath}
\usepackage{graphicx}

\title{Tempo：免选举的租约定序共识协议（插件演示样例）}
\author{张三}

\begin{document}
\maketitle

\begin{abstract}
工程上广泛使用的共识协议，大多把「请求排序」与「状态持久化」这两件事都压在同一个稳定领导者身上：
领导者一旦失效，所有写入都要等新一轮选举结束才能恢复。
本文提出 Tempo——一种把排序权与持久化权分离的共识协议。
排序权由按纪元（epoch）轮转的仲裁节点（arbiter）以短租约（lease）的形式授予，
被接受的请求会拿到一枚定序令牌（Ordering Token, OT）；持久化仍然由多数派写入保证。
仲裁节点的身份是纪元号的确定性函数，不需要通过选举产生，
因此换主只需要一轮仲裁，不存在选举停摆期。
在部分同步网络与崩溃失效模型下，我们给出令牌唯一性与纪元单调性两条引理，
并据此证明 Tempo 满足线性一致性。
\end{abstract}

\section{引言}
共识协议要解决的问题是：让一组可能失效的副本对同一个操作序列达成一致。
Paxos~\cite{lamport1998} 与 Raft~\cite{ongaro2014} 都把「先选出稳定领导者」当作前提，
由领导者为请求排序、副本按序追加日志~\cite{lamport2001}。
这个设计把复杂性收敛到了单一节点，代价是可用性也绑在了这个节点上：
领导者失效期间，写入必须停摆到新一轮选举结束。
下面这一行故意引用了一个不存在的 label，用于演示引用校验功能：
参见第~\ref{sec:intro}~节的动机说明。
% TODO: 补充近三年（2023--2025）无领导者共识方向的综述文献

本文的贡献可以概括为三点：
\begin{itemize}
  \item \textbf{排序与持久化解耦}：定序由仲裁节点单独授权，写入由多数派确认，两者互不阻塞；
  \item \textbf{免选举换主}：仲裁节点的身份是纪元号的确定性函数，换主不需要任何投票；
  \item \textbf{可分析的快路径}：正常路径 1 RTT，竞争路径 2 RTT，崩溃恢复只做尾部对齐。
\end{itemize}

\section{相关工作与设计动机}\label{sec:related}
租约（lease）最早由 Gray 与 Cheriton~\cite{gray1989} 提出，
用于在分布式文件缓存中避免反复确认；它提供了时间维度上的互斥，但本身不提供共识。
Paxos 之后的工程实现（Multi-Paxos、ZAB、Raft）选择了「稳定领导者 + 日志复制」的路线，
把一轮协商摊薄成常态化的领导权。
EPaxos~\cite{moraru2013} 走向另一个极端：它不设领导者，用依赖集描述命令之间的偏序，
冲突稀疏时可以单轮提交，冲突密集时则需要额外的提交轮次。
Tempo 借用了「不过度依赖单点」这个动机，但换了一条实现路径：
把临时排序权显式化成一枚租约，并让租约携带由多数派确认的纪元栅栏，
从而让「排序」这件事可以直接参与安全性论证。

\section{系统模型与假设}\label{sec:model}
我们考虑 $n = 2f + 1$ 个副本组成的集群，任意时刻最多 $f$ 个副本失效。
失效模型是崩溃失效（crash-stop）：失效节点不再产生消息，也不做拜占庭行为。
网络是部分同步的：存在一个未知时刻 GST，此后任意消息的传输延迟不超过已知上界 $\Delta$。
多数派（quorum）定义为任意 $\lceil (n+1)/2 \rceil$ 个副本。以下三点假设贯穿全文：
\begin{enumerate}
  \item 副本的本地时钟只要求单调不回退，不要求彼此同步；
  \item 客户端允许重试，因此协议只需要保证「至多生效一次」的语义；
  \item 消息可能丢失、重复与乱序，但不会损坏，也不会被篡改。
\end{enumerate}

\section{协议设计}\label{sec:method}
Tempo 把一次写入拆成两个可以独立推进的阶段：先向仲裁节点取定序令牌，
再把请求按令牌顺序写入多数派。图~\ref{fig:flow} 给出了正常路径上的消息流。
\begin{figure}[htbp]
  \centering
  % \includegraphics[width=.85\linewidth]{tempo-flow.pdf}
  \caption{正常路径的消息流：客户端 $\to$ 提案节点 $\to$ 仲裁节点 $\to$ 多数派}\label{fig:flow}
\end{figure}

\subsection{混合逻辑时钟与纪元栅栏}\label{sec:hlc}
副本维护一个混合逻辑时钟（Hybrid Logical Clock, HLC）~\cite{lamport1978}，
它取本地物理时间与收到的最大逻辑时间中的较大者，只要求单调不回退。
纪元（epoch）是排序权的时间切片。当仲裁节点需要开启新纪元 $e+1$ 时，
它先向 $n-1$ 个副本发出探测，收集其中 $\lceil (n+1)/2 \rceil - 1$ 个确认，
再按下式推进基准时钟：
\begin{equation}\label{eq:fence}
  t_{e+1} = \max\left(t_{e},\ h_{\max}\right) + 1
\end{equation}
其中 $h_{\max}$ 是确认集合 $\mathcal{A}$ 中各副本报告的 HLC 的最大值。
式~\eqref{eq:fence} 得到的 $t_{e+1}$ 称为纪元栅栏：
任何属于纪元 $e+1$ 的请求，其时间戳都不允许小于这道栅栏。

\subsection{定序令牌与租约}
仲裁节点在纪元 $e$ 内为本地的序号 $s$ 单调递增地签发定序令牌：
\begin{equation}\label{eq:token}
  \mathrm{OT} = \langle e,\ s \rangle, \qquad
  \langle e,\ s \rangle < \langle e',\ s' \rangle \Leftrightarrow
  e < e' \lor (e = e' \land s < s')
\end{equation}
令牌之间按字典序（lexicographic order）比较，因此任意两枚令牌都可以比出大小，
这也让提交序列在副本之间有了唯一确定的排列。
一枚令牌只在租约期内有效：租约时长 $L$ 需要在「仲裁节点的负载」与「失效检测的迟缓」之间折中，
工程上通常取 $L$ 为几倍的消息往返时间。\footnote{本文取 $L = 3\Delta$，并在实现中把它做成可配置项。}
令牌有效当且仅当：
\begin{equation}\label{eq:lease}
  \mathrm{valid}(\mathrm{OT}) \Leftrightarrow
  \mathrm{now} < \mathrm{expire}(\mathrm{OT}) \land
  \mathrm{Arb}(e) = a \land
  |\mathrm{Ack}_{e}| \geq \lceil (n+1)/2 \rceil
\end{equation}
其中 $\mathrm{Arb}(e)$ 是由纪元号唯一确定的仲裁节点，
$\mathrm{Ack}_{e}$ 是纪元 $e$ 收集到的确认集合。

\subsection{正常路径与竞争路径}
正常路径上，提案节点把请求连同租约一起广播给多数派，收到多数派的写入确认后即可提交，
总耗时为 1 RTT。若租约已经过期，或者仲裁节点在探测阶段没能在 $L$ 内收到足够确认，
请求就转入竞争路径：提案节点重新收集 HLC、推进纪元栅栏，再取一枚新令牌，总耗时为 2 RTT。
两条路径的差别只体现在延迟上，不影响安全性。

\section{正确性论证}\label{sec:correct}
本节给出 Tempo 的安全性论证。记号约定：$\mathrm{OT}_{x}$ 表示请求 $x$ 的定序令牌，
$\mathrm{Ack}_{e}$ 表示纪元 $e$ 中收集到确认的副本集合，$t_{e}$ 表示纪元 $e$ 的基准时钟。

\textbf{引理 1（令牌唯一性）} 同一纪元内，仲裁节点不会把同一个序号签发给两个不同的请求。
\emph{证明}：仲裁节点对序号 $s$ 的分配是本地单调递增的，而且签发之前必须通过租约有效性检查~\eqref{eq:lease}；
序号一旦被签发，单调性使得后续请求只能拿到更大的序号。证毕。

\textbf{引理 2（纪元单调性）} 若纪元 $e' > e$，则 $t_{e'} > t_{e}$。
\emph{证明}：推进纪元需要收集 $\lceil (n+1)/2 \rceil - 1$ 个副本的确认，并按式~\eqref{eq:fence} 取最大 HLC 再加一。
两次推进的确认集合大小之和超过 $n$，因此必存在一个副本同时出现在两次确认集合里；
该副本的 HLC 单调不回退，于是后一次推进读到的最大值不小于前一次，即 $t_{e'} > t_{e}$。证毕。

\textbf{定理 1（线性一致性）} 在部分同步网络中，Tempo 的提交序列等价于某个顺序执行。
\emph{证明}：由引理 1，同一纪元内不存在两枚相同的令牌；由引理 2，纪元之间不会出现倒序。
于是全部提交都可以按式~\eqref{eq:token} 的字典序排列成一个唯一的全序。
每次提交都写入了多数派，任意两次提交的写入集合相交，交集中的副本按同一顺序追加请求，
因此任一时刻任一读取看到的都是这个全序的前缀。证毕。

\section{对比与评估}\label{sec:eval}
表~\ref{tab:cmp} 给出四类共识协议在 5 节点集群上的开销对比。
\begin{table}[htbp]
  \centering
  \caption{四类共识协议在 5 节点集群上的开销对比}\label{tab:cmp}
  \small
  \begin{tabular}{lcccc}
    \hline
    协议 & 需领导者 & 快路径 RTT & 换主 & 恢复 \\
    \hline
    Multi-Paxos & 是 & 1 & 选举·随机退避 & 补齐空洞 \\
    Raft & 是 & 1 & 选举·随机超时 & 截断冲突 \\
    EPaxos & 否 & 1--2 & 无需换主 & 依赖集重放 \\
    Tempo（本文） & \textbf{否} & \textbf{1} & \textbf{一轮仲裁} & \textbf{尾部对齐} \\
    \hline
  \end{tabular}
\end{table}

Multi-Paxos 与 Raft 在领导者稳定时都能达到 1 RTT，但它们把可用性绑定在领导者身上：
领导者失效之后，集群要等到选举超时（通常是 $150\sim 300$ ms）才能恢复服务。
Tempo 的仲裁轮只需要一次往返，而且在 $\Delta$ 有界的假设下是可以预期的：
换主不需要任何节点就「谁是新仲裁者」达成共识，因为新仲裁者由纪元号直接算出。
EPaxos 同样不需要领导者，但它的事务提交流程与命令之间的依赖关系有关，冲突密集时延迟会上升；
Tempo 把排序权收敛到单枚租约上，因此路径长度与冲突密度无关。
代价是仲裁节点在租约期内会成为热点。

\section{局限与未来工作}
本文只讨论了单组（single-shard）场景。跨组事务需要额外引入协调者，
这会重新带回 Tempo 试图消除的那个单点。
另一处局限是纪元栅栏依赖时钟的单调性：如果某个副本发生时钟回退，栅栏必须显式拒绝它的确认。
未来工作包括 % FIXME: 补写跨组事务的形式化模型与时钟回退的处理
把租约时长做成自适应参数，并在更大规模的集群上补做实验。

\section{结论}
本文提出了 Tempo：一种把排序权与持久化权分离的共识协议。
它用一枚由多数派确认过的定序令牌替代了领导者选举，使换主成为一次有界的仲裁轮。
我们给出了令牌唯一性与纪元单调性两条引理，并据此证明了线性一致性。
在正常路径上，Tempo 与 Raft 具有相同的往返次数，但换主过程不再需要选举停摆期。

\begin{thebibliography}{9}
\bibitem{lamport1978} Lamport L. Time, clocks, and the ordering of events in a distributed system. Communications of the ACM, 1978.
\bibitem{gray1989} Gray C, Cheriton D. Leases: an efficient fault-tolerant mechanism for distributed file cache consistency. SOSP, 1989.
\bibitem{lamport1998} Lamport L. The part-time parliament. ACM Transactions on Computer Systems, 1998.
\bibitem{lamport2001} Lamport L. Paxos made simple. ACM SIGACT News, 2001.
\bibitem{moraru2013} Moraru I, Andersen D G, Kaminsky M. There is more consensus in egalitarian parliaments. SOSP, 2013.
\bibitem{ongaro2014} Ongaro D, Ousterhout J. In search of an understandable consensus algorithm. USENIX ATC, 2014.
\end{thebibliography}

\end{document}
